\section{Morfismos planos de preesquemas localmente noetherianos}

Sean $X$, $Y$ dos preesquemas localmente noetherianos, y $f:X\to Y$ un
morfismo. Para todo $y\in Y$, la fibra
$f^{-1}(y)=X\times_Y\operatorname{Spec}(\boldsymbol{k}(y))$ es también un preesquema
localmente noetheriano, pues basta verlo cuando
$Y=\operatorname{Spec}(A)$ y $X=\operatorname{Spec}(B)$ son espectros de
anillos noetherianos, y entonces $B\otimes_A \boldsymbol{k}(y)$ es un anillo de fracciones
del anillo cociente $B/\mathfrak j_yB$, por tanto es noetheriano. Nos
proponemos, en primer lugar, poner en relación en este parágrafo las propiedades
de $X$, de $Y$ y de las fibras $f^{-1}(y)$ (cuando $y$ recorre $Y$), bajo la
hipótesis de que el morfismo $f$ es \emph{plano}. Las cuestiones tratadas se reducirán
al estudio de las relaciones entre los anillos locales noetherianos
$\mathcal O_y$, $\mathcal O_x$ y
$\mathcal O_x\otimes_{\mathcal O_y}\boldsymbol{k}(y)$, siendo el homomorfismo
$\mathcal O_y\to\mathcal O_x$ local y haciendo de $\mathcal O_x$ un
$\mathcal O_y$-módulo plano. En los n\textsuperscript{os} 6.11 a 6.13
(bastante separados del resto del parágrafo por su naturaleza «absoluta» más
que «relativa»), aplicamos algunos de los resultados precedentes para obtener
unos

\oldpage[IV-2]{135}

criterios que permitan afirmar que el \emph{lugar singular} (o ciertos
conjuntos análogos) de ciertos preesquemas son conjuntos cerrados; estos
criterios desempeñarán un papel importante en el §~7.
